\section{Conclusion}
\label{sec:conclusion}
We introduced the \textbf{Red Queen G\"odel Machine (\ourmethod)}, a framework in which evaluators co-evolve with the agents they score. Our results suggest two broad principles. First, co-evolving the evaluator alongside the generator enables improvement on hard-to-verify tasks such as paper writing and proof writing, where a fixed benchmark cannot score the artifact directly, mirroring co-evolutionary dynamics observed in nature~\citep{van1973new}. Second, co-evolved systems match or exceed fixed-evaluator baselines while often being more token-efficient, which we posit is due to shared expansion costs, more heterogeneous utility signals, and a curriculum-like effect in which a progressively stricter evaluator hardens the population over time. Underlying both is the framework's accommodation of non-stationary utilities: because controlled utility evolution lets the search objective change at epoch boundaries while keeping per-epoch guarantees intact, interventions a fixed objective cannot express, such as the adversarial term that corrects the reviewer's self-preference bias, become available mid-search.

\subsection{Future Research Directions}
\label{sec:future}
Our investigation is preliminary, drawn from short search horizons we intend to extend. The framework's premise is that longer co-evolution compounds these effects, extending the curriculum so that each generation of agents faces a stricter and more discerning judge. The limit of this improvement is governed by the anchor. Anchoring every replacement to fixed ground truth keeps evaluators accurate but confines them to the anchor's own decision boundary. Improvement far beyond the benchmark must therefore come from the objectives layered on top of the ground-truth signal, with the anchor serving as a guardrail against drift. The adversarial reviewer is the first such instance: it roughly maintains accuracy on the ground truth while the adversarial term discovers a human--machine decision boundary the benchmark never specified. Richer objectives of this kind are, we believe, the path toward systems that bootstrap their own evaluation beyond the reach of static benchmarks. We identify two limitations with the current framework. First, evaluator quality is only as good as its anchor: a weak or biased anchor could yield uninformative evaluators. Second, our theoretical guarantees are epoch-local, covering improvement only within an epoch. We discuss the limitations of \theourmethod in full in \cref{app:limitations}, and will update this work as we explore these future research directions.